% ----------------------------------------------------------
\begin{frame}{Views -- What and Why}
  A \textbf{VIEW} is a named, stored \texttt{SELECT} query that behaves like a virtual table.
  The underlying data is not duplicated -- the view runs its query fresh on every access.

  \begin{block}{Why Use Views?}
    \begin{itemize}
      \item \textbf{Simplify complex queries}: wrap a multi-table JOIN into a single named view;
            application code just selects from the view.
      \item \textbf{Access control}: expose only specific columns/rows to certain users without
            granting access to the underlying tables.
      \item \textbf{Logical abstraction}: rename or reformat columns without changing the schema.
      \item \textbf{Backward compatibility}: when a table is restructured, a view with the old
            column names can shield existing applications from breaking changes.
    \end{itemize}
  \end{block}

  \begin{alertblock}{Views are Not Materialized (by default)}
    A standard view re-executes its query every time it is accessed. For performance-critical
    use, some DBMS (PostgreSQL, Oracle) offer \textbf{materialized views} that cache results.
    MySQL does not natively support materialized views.
  \end{alertblock}
\end{frame}

% ----------------------------------------------------------
\begin{frame}[fragile]{Creating and Using Views}
  \begin{lstlisting}[style=sqlstyle]
-- Create a view: all tracks with artist and album info
CREATE VIEW v_track_details AS
SELECT t.title   AS track_title,
       t.duration_sec,
       al.title  AS album_title,
       al.year,
       ar.name   AS artist_name
FROM   Track  AS t
JOIN   Album  AS al ON t.album_id   = al.id
JOIN   Artist AS ar ON al.artist_id = ar.id;

-- Use the view like a regular table
SELECT * FROM v_track_details WHERE artist_name = 'Aurora';

-- Remove the view
DROP VIEW IF EXISTS v_track_details;
  \end{lstlisting}

  \begin{block}{Updatable Views}
    Simple views (single table, no GROUP BY, no DISTINCT) can be used with
    \texttt{INSERT} / \texttt{UPDATE} / \texttt{DELETE}. Complex views with JOINs
    or aggregations are read-only.
  \end{block}
\end{frame}

% ----------------------------------------------------------
\begin{frame}[fragile]{Triggers -- Concept and Syntax}
  A \textbf{TRIGGER} is a stored procedure that the DBMS executes \emph{automatically}
  in response to an \texttt{INSERT}, \texttt{UPDATE}, or \texttt{DELETE} on a specific table.

  \begin{lstlisting}[style=sqlstyle]
CREATE TRIGGER trigger_name
  { BEFORE | AFTER }
  { INSERT | UPDATE | DELETE }
  ON table_name
  FOR EACH ROW
BEGIN
  -- SQL statements using NEW.col and/or OLD.col
END;
  \end{lstlisting}

  \begin{block}{NEW and OLD}
    \begin{itemize}
      \item \texttt{NEW.column} -- the new value being inserted or the updated value.
      \item \texttt{OLD.column} -- the value that existed before the update or deletion.
      \item \texttt{BEFORE INSERT}: only \texttt{NEW} available.
      \item \texttt{BEFORE UPDATE}: both \texttt{OLD} and \texttt{NEW} available.
      \item \texttt{AFTER DELETE}: only \texttt{OLD} available.
    \end{itemize}
  \end{block}
\end{frame}

% ----------------------------------------------------------
\begin{frame}[fragile]{Trigger Example -- BEFORE INSERT}
  \textbf{Requirement:} Automatically set a track's \texttt{created\_at} timestamp and
  normalise its title to trimmed title case before inserting.

  \begin{lstlisting}[style=sqlstyle]
DELIMITER //
CREATE TRIGGER trg_track_before_insert
BEFORE INSERT ON Track
FOR EACH ROW
BEGIN
    -- Ensure duration is never negative
    IF NEW.duration_sec < 0 THEN
        SET NEW.duration_sec = 0;
    END IF;

    -- Trim whitespace from the title automatically
    SET NEW.title = TRIM(NEW.title);
END //
DELIMITER ;

-- Test it:
INSERT INTO Track (title, duration_sec, album_id)
VALUES ('  Running With the Wolves  ', -5, 1);
-- Stored as: title='Running With the Wolves', duration_sec=0
  \end{lstlisting}
\end{frame}

% ----------------------------------------------------------
\begin{frame}[fragile]{Trigger Example -- AFTER INSERT (Audit Log)}
  \textbf{Requirement:} Record every new track insertion in an audit table.

  \begin{lstlisting}[style=sqlstyle]
-- Create an audit table first
CREATE TABLE Track_Audit (
    id         INT AUTO_INCREMENT PRIMARY KEY,
    track_id   INT,
    action     VARCHAR(10),
    changed_at TIMESTAMP DEFAULT CURRENT_TIMESTAMP
);

DELIMITER //
CREATE TRIGGER trg_track_after_insert
AFTER INSERT ON Track
FOR EACH ROW
BEGIN
    INSERT INTO Track_Audit (track_id, action)
    VALUES (NEW.id, 'INSERT');
END //
DELIMITER ;
  \end{lstlisting}

  \begin{block}{BEFORE vs AFTER}
    Use \texttt{BEFORE} to validate or modify values before they are saved.
    Use \texttt{AFTER} for side effects (audit logs, cascading updates to other tables).
  \end{block}
\end{frame}

% ----------------------------------------------------------
\begin{frame}[fragile]{Stored Procedures}
  A \textbf{STORED PROCEDURE} is a named, reusable block of SQL stored in the database.
  It can accept parameters, contain control-flow logic, and return result sets.

  \begin{lstlisting}[style=sqlstyle]
DELIMITER //
CREATE PROCEDURE add_track(
    IN  p_title        VARCHAR(200),
    IN  p_duration_sec INT,
    IN  p_album_id     INT,
    OUT p_new_id       INT
)
BEGIN
    INSERT INTO Track (title, duration_sec, album_id)
    VALUES (p_title, p_duration_sec, p_album_id);

    SET p_new_id = LAST_INSERT_ID();
END //
DELIMITER ;

-- Call the procedure
CALL add_track('Runaway', 237, 3, @new_track_id);
SELECT @new_track_id;  -- retrieve the OUT parameter
  \end{lstlisting}

  \begin{itemize}
    \item Reduces network round-trips: logic runs server-side.
    \item \texttt{IN}: input parameter \quad \texttt{OUT}: output parameter \quad \texttt{INOUT}: both.
  \end{itemize}
\end{frame}

% ----------------------------------------------------------
\begin{frame}{Views, Triggers, and Procedures -- Comparison}
  \begin{center}
  \renewcommand{\arraystretch}{1.35}
  \begin{tabularx}{\linewidth}{@{}lXX@{}}
    \toprule
    \textbf{Object} & \textbf{Purpose} & \textbf{When to Use} \\
    \midrule
    \textbf{VIEW}
      & Virtual table wrapping a SELECT query
      & Simplify queries; restrict column access; logical abstraction \\[4pt]
    \textbf{TRIGGER}
      & Auto-execute SQL on INSERT/UPDATE/DELETE
      & Enforce business rules; maintain audit logs; cascade custom logic \\[4pt]
    \textbf{STORED PROCEDURE}
      & Reusable named SQL block with parameters
      & Encapsulate multi-step operations; reduce client-server round-trips \\
    \bottomrule
  \end{tabularx}
  \end{center}

  \begin{alertblock}{Caution with Triggers}
    Triggers execute invisibly -- a DML statement may have hidden side effects that are hard to
    debug. Document triggers clearly, and avoid complex business logic that belongs in the
    application layer.
  \end{alertblock}
\end{frame}
